% =============================================================================
% 04_elastic_net.tex — Problem 4: Elastic Net and Neural-Feature Bridge
%                      (P5 — Tài)
% =============================================================================
\section{Elastic Net and Neural-Feature Bridge}
\label{sec:elastic-net}

\subsection{Research Question and Source Map}
\label{sec:research-question}

\textbf{Research question.}  When does an $\ell_2$ penalty on neural-network
weights behave like classical Ridge regression, and why must that connection be
qualified for adaptive optimisation and distinguished from dropout?  The first
source below is a primary research paper not cited in the project brief.

\begin{table}[H]
\centering
\small
\caption{Source map linking methodological claims to primary evidence.}
\label{tab:source-map}
\begin{tabularx}{\textwidth}{>{\raggedright\arraybackslash}X
  >{\raggedright\arraybackslash}p{0.22\textwidth}
  >{\raggedright\arraybackslash}X
  >{\raggedright\arraybackslash}X}
\toprule
\textbf{Claim} & \textbf{Exact Source and Locator} &
\textbf{What It Establishes} & \textbf{What It Does Not Establish} \\
\midrule
A quadratic weight penalty yields a decay term in gradient learning. &
\citet{Krogh1991}, Section 1, Equations (1)--(2). &
Adds $\lambda\sum_iw_i^2/2$ to a neural-network loss and derives the
$-\lambda w_i$ update term. &
Does not analyse Adam or prove that every modern ``weight decay''
implementation is equivalent. \\[4pt]
$\ell_2$ regularisation and weight decay separate under adaptive updates. &
\citet{LoshchilovHutter2019}, abstract and Section 2. &
Shows equivalence for standard SGD after learning-rate rescaling and
inequivalence for adaptive gradient methods; motivates decoupled AdamW. &
Does not make the present frozen-feature experiment a trained neural network. \\[4pt]
Dropout randomly masks units during training and uses the unthinned network at
test time. &
\citet{Srivastava2014}, Section 2 and Figure 2. &
Defines stochastic activation masking and the corresponding test-time scaling. &
Does not establish permanent sparsity or deterministic pruning of learned
weights. \\
\bottomrule
\end{tabularx}
\end{table}

\subsection{Elastic Net on Original Predictors}
\label{sec:enet-original}

The Elastic Net \citep{ZouHastie2005} combines $\ell_1$ and $\ell_2$ penalties:
\[
 \hat{\bbeta}^{\mathrm{ENet}}
 =\argmin_{\bbeta}\left\{
 \frac{1}{2n}\norm{\by-\bX\bbeta}_2^2
 +\lambda\left[\alpha\norm{\bbeta}_1
 +\frac{1-\alpha}{2}\norm{\bbeta}_2^2\right]\right\}.
\]
Using the shared folds, $\lambda$ was tuned separately for
$\alpha\in\{0.10,0.25,0.50,0.75,0.90\}$.  The pair with the smallest CV MSE
was $\alpha=0.90$, $\lambda_{\min}=0.1061$, with CV RMSE 4.242 and nine
nonzero slopes.

\includefigure{0.75\textwidth}{fig_p4_enet_cv}%
  {Elastic Net shared-fold CV performance across candidate $\alpha$ values.}%
  {fig:enet-alpha}

\includefigure{0.75\textwidth}{fig_p4_enet_path}%
  {Elastic Net coefficient path for the selected $\alpha=0.90$.}%
  {fig:enet-path}

\includefigure{0.72\textwidth}{fig_p4_enet_coef}%
  {Elastic Net coefficients at the selected $(\alpha,\lambda)$.}%
  {fig:enet-coef}

\input{../output/tables/tab_p4_enet_alpha}
\input{../output/tables/tab_p4_enet_coef}

The selected $\alpha$ lies close to Lasso, and its active set and CV error are
accordingly similar.  Its nonzero $\ell_2$ component still shares shrinkage
across correlated predictors, while the dominant $\ell_1$ component preserves
sparsity; it therefore lies between Ridge's dense solution and Lasso's exact
thresholding.

\subsection{Fixed ReLU Features}
\label{sec:relu-features}

With random-feature seed \texttt{240401}, the standardised predictors are
mapped once to $M=200$ features
\[
 h_m(\bx)=\max\{\bm a_m^\top\bx+c_m,0\},\qquad
 a_{jm}\sim N(0,1/p),\quad c_m\sim N(0,0.5^2).
\]
The same sampled matrix $A$ and bias vector $\bm c$ are applied to training and
test predictors.  One constant hidden feature is detected from $H_{\mathrm{train}}$
alone and removed from both matrices; centring and scaling of the remaining
hidden features are also learned on training rows only.  The random hidden
weights and biases are fixed.  Only the output intercept and slopes are fitted,
with $\lambda$ (and the preselected candidate $\alpha$ values) chosen by the
shared training folds.

\includefigure{0.72\textwidth}{fig_p4_neural_ridge_cv}%
  {Shared-fold CV curve for Ridge on the fixed ReLU features.}%
  {fig:neural-ridge-cv}

\includefigure{0.72\textwidth}{fig_p4_neural_lasso_cv}%
  {Shared-fold CV curve for Lasso on the fixed ReLU features.}%
  {fig:neural-lasso-cv}

\includefigure{0.72\textwidth}{fig_p4_neural_active}%
  {Active output weights for the three fixed-feature models.}%
  {fig:neural-active}

\input{../output/tables/tab_p4_neural_summary}

The nonlinear expansion does not improve shared-fold prediction.  Neural Ridge
has CV RMSE 4.784 with all 199 usable features active; Neural Lasso and Neural
Elastic Net reduce this to about 4.390 with 31 active features, but remain worse
than the original-predictor Lasso and Elastic Net (about 4.242).  With only 201
training cases, the 199-dimensional random representation adds variance and
depends on a single frozen draw rather than learning task-specific features.

\subsection{Locked Declaration and Final Holdout}
\label{sec:holdout}

\noindent
\fbox{\parbox{0.95\textwidth}{%
\textbf{Locked Analysis Declaration (before any \texttt{y\_test} metric)}\\[4pt]
Training-only centring/scaling is applied unchanged to test predictors.
Candidate specifications are OLS; Ridge at $\lambda_{\min}=0.627974$; Lasso at
$\lambda_{\min}=0.095447$; Elastic Net at $\alpha=0.90$ and
$\lambda_{\min}=0.106052$; and the fixed-feature models using seed 240401 and
$M=200$.  Based on training CV, the deployment model is the
\textbf{original-predictor Lasso}.  The locked primary metrics are RMSE and MAE,
with $R^2$ supplementary.}}

\input{../output/tables/tab_p4_holdout}

\includefigure{0.72\textwidth}{fig_p4_holdout_compare}%
  {RMSE and MAE for all candidates in the one final holdout evaluation.}%
  {fig:holdout-comparison}

\includefigure{0.62\textwidth}{fig_p4_actual_vs_pred}%
  {Actual versus predicted \texttt{brozek} for the predeclared Lasso on the holdout set.}%
  {fig:actual-vs-predicted}

The holdout evidence supports the declaration: Lasso ranks first with RMSE
4.251, MAE 3.524, and $R^2=0.683$.  Elastic Net is close (RMSE 4.274), but no
model was retuned or replaced after viewing these results.  The frozen ReLU
models rank below every original-predictor candidate; their best holdout RMSE
is 4.757 for Neural Ridge.

\subsection{Deep-Learning Connection and Limitations}
\label{sec:dl-connection}

For a fixed matrix $H$, an $\ell_2$ penalty on the learned output slopes is
exactly Ridge regression.  In an adaptively optimised neural network, however,
adding an $\ell_2$ term to the loss need not equal multiplicatively decaying the
weights; the update geometry depends on the optimiser
\citep{LoshchilovHutter2019}.  A zero output slope here removes one fixed hidden
feature from this scalar predictor, but it does not make the input-to-hidden
matrix $A$ sparse.  Dropout instead samples temporary activation masks during
training and uses a scaled unmasked network at test time
\citep{Srivastava2014}; this is not permanent pruning.  Most importantly, our
single hidden layer is frozen, so the experiment does not establish a result
about end-to-end training or a fully trained deep network.

The experiment has at least three limitations.  First, it uses only one random
feature seed, so the apparent ranking may change with another draw.  Second,
$n_{\mathrm{train}}=201$ is small relative to 200 proposed hidden features,
making a flexible nonlinear representation difficult to estimate reliably.
Third, the hidden weights are not learned and only one width and bias scale are
studied.  These restrictions make the result a controlled regularisation
comparison, not a general assessment of neural networks.
